\documentclass[11pt]{article}
\usepackage[margin=27mm]{geometry}
\usepackage{amsmath,amssymb,fontspec,unicode-math,booktabs,array,hyperref}
\setmainfont{STIXGeneral}
\setmathfont{STIX Two Math}
\setmonofont{Menlo}
\hypersetup{colorlinks=true,urlcolor=blue}
\setlength{\parindent}{0pt}
\setlength{\parskip}{5pt}
\title{Affine proof dependencies}
\author{\texttt{human\_compression}}
\date{Logical graph candidate; independent semantic review pending}
\begin{document}
\maketitle
\vspace{-1.2em}

The admitted contract recovers the main results of Hou--Feng--Leemans,
\emph{Regular 3-polytopes of order $2^np$}, in the bound 17-page preprint
arXiv:2001.02945v2 (14 July 2024). Its 33 records include principal results,
borrowed inputs, explicit exclusions and a source correction. They are not
33 independent new theorems.

The paired graph compares the dependencies actually used in that source
with those of the complete current common paper and application. Arrows point
from a prerequisite to its use. Historical arrows document the source's
argument; they do not retroactively validate an incorrect intermediate claim.
Attribution, formal evidence, comparisons and operational links are kept out
of mathematical prerequisite unions.

\textbf{The common mechanisms.}
The kernel-profile criterion reduces marked intersections to intersections
inside an explicitly identified kernel. Supported coordinate kernels give
the odd-factor constructions. For the two cubic constructions, one exact
abstract kernel presentation and a faithful affine model identify the whole
infinite group. Invariant translation periods then give the literal finite
presentations; finite overlap differences prove IP for every multiplier.

\begin{tabular}{@{}>{\raggedright\arraybackslash}p{.25\linewidth}>{\raggedright\arraybackslash}p{.69\linewidth}@{}}
\toprule
Current conclusion & Actual principal route \\
\midrule
Odd-prime type restriction & Burnside and the dihedral quotient by $O_2$, then the adjacent-label argument. \\
Two odd edge factors & Complete square presentation and periods; supported kernels; two cyclic layers give the reverse presentation bound. \\
Three $\{6,4\}$ families & Exact index-48 kernel, normalized affine transitions and all source-word lattices; period and overlap criteria. \\
Full $\{6,2^s\}$ range & One-edge extension, exceptional 384 seed and exact source-period adapter; cyclic-core descent and complete catalogue exclusion. \\
Three $\{6,6\}$ families & Exact index-48 kernel; invariant closure of raw transitions; all source periods and the nonzero overlap obstruction. The order-96 case has an explicit tetrahedral construction. \\
\bottomrule
\end{tabular}

The one-edge extension has its own proof requiring only one deletion. Its
relation to the general supported-kernel theorem is recorded as a comparison,
not an imported proof. The complete common paper remains the declared
boundary, including its unused deletion corollary and smaller-period lemma.

\newpage
\textbf{The historical routes and their qualifications.}
The original cubic proofs separately rewrite three finite-index subgroups,
recognize three finite base geometries, take characteristic power subgroups,
and apply the quotient criterion. The graph records the source's Magma,
Reidemeister--Schreier, atlas and classical inputs at those uses. The current
proof replaces those particular routes with the explicit common mechanism
on the previous page; it still pays for its actual exact computations.

In the bound preprint, Lemma 4.5's first $R_0$ equality has the wrong sign.
The source's Theorem 1.3 route invokes that lemma. The current paper proves
the obstruction, supplies the negative exponent and identifies the corrected
family. These are separate graph records. The graph also flags Corollary
5.2's reference to Theorem 1.2: its displayed family argument uses Theorem
5.1. These qualifications concern the bound preprint, not an unexamined
journal edition.

\textbf{Programs and retained inputs.}
Six whole mathematical programs are bound: the two exact abstract kernel
calculations, the two integer affine checks, the exceptional-seed/sign
calculation, and the complete order-384 catalogue search. Embedded controls
remain included. The two portable drivers and the separate finite-control
file are auxiliary evidence, not mathematical premises. A file-hash link
between the cubic programs is operational provenance, not a claim that one
executes the other.

Burnside's theorem, Lucchini's cyclic-core bound, the regular-polytope
correspondence, exact presentation algorithms, GAP and SmallGroups remain
explicit retained inputs. The current rank-three intersection, square-map
presentation, tight case and order bound have local proofs. Historical atlas
and square-classification citations are not silently imported into those
proofs. External proof and graphic costs remain unknown, not zero.

\textbf{Formal evidence and the closed contract.}
The accepted nine Lean declarations cover two smaller-period statements,
three sign-obstruction identities and four quotient-intersection criteria
used by the formal period proof. They do not certify the affine coordinates,
actual subgroup presentations, catalogue completeness, or source-specific
geometry. No fresh Lean or group census is inferred from checking this graph.

Problem 1.1, the separate Table 1 census and uniqueness of the two 384 seeds
remain excluded. The final Type (2) existence question was already answered
historically; the current application includes a complete local binary-module
construction with explicit attribution. It earns no new question credit.

\textbf{What this packet establishes.}
All 33 accepted records and 17 current proof contexts are bound. The checker
verifies typed acyclic dependencies, prerequisite unions, exact artifact
hashes and deliberate corruption controls. Independent semantic review must
still confirm that each arrow describes an actual proof use. Context spans
may overlap: this is not a text partition. Node and edge counts are only
inventories; no weight, compression factor or new result credit is assigned.
\end{document}
